\cotefile{156-9}
\scanpage{50}
\mspage{46}
% archive sheet number: 49
On trouve donc que la structure d'\uline{atelier spatial} peut se décrire par les données
\[
(\unsure{\mathcal A|\mathscr A},\trianglelefteq,\mathrel{|\unsure{o|\circ}|})
\]
avec les seuls axiomes

\uline{At spat 1}\quad a) $\trianglelefteq$ est une relation d'ordre.

\hspace*{4em}b) $\mathrel{|o|}$ est symétrique \unsure{antiréflexive}

\uline{At spat 2}\quad Si $X\triangleleft Y$, on a $X\mathrel{|o|}Y$.

\uline{At spat 3}\quad L'application
\[
\begin{aligned}
\mathcal A &\longrightarrow \deleted{\unsure{\mathfrak S|S}}\,\mathrm{Fig}\,\text{él}\,(\unsure{P|\mathcal P}(\mathcal A))\\
X &\longmapsto \text{Dép}\,X
 = \bigl(\{X'^{\unsure{\circ|o}}\mid X'\unsure{\mathrel{\ooalign{\hfil$\trianglelefteq$\hfil\cr\hfil\raisebox{0.32ex}{\scriptsize$-$}\hfil\cr}}|\trianglelefteq}X\},\,\deleted{\text{\unsure{avec}}\,\unsure{\leq}}\,+\\
 &\hspace{2.6em}\deleted{\text{où }X'^{\circ}=\{Y}\quad
 \inserted{\begin{aligned}
 &\text{relation d'ordre sur Dép }\unsure{X|X'}\\[-2pt]
 &\text{déduite de celle de}\\[-2pt]
 &\widetilde X=\{x'\in\mathcal A\mid x\trianglelefteq X\}
 \end{aligned}}\bigr)
\end{aligned}
\]
est \uline{injective}. (\unsure{On|on} a \unsure{a} posé, pour $Y\in\mathcal A$,
\[
Y^{\circ}=\{Z\in\mathcal A\mid\forall Z'\in\mathcal A,\ Z'\mathrel{|o|}Y\Rightarrow Z'\mathrel{|o|}Z\}
\]

Je vais essayer de \unsure{reformuler} la structure, avec $\Sigma=\Sigma_{\mathcal A}$ comme ensemble de base.
\margin{$(\Sigma,$}

\uline{Définition} On appelle \uline{\unsure{spatios|spatio}} (ensembliste) un \deleted{\unsure{ensemble ordonné}}\inserted{triple} $(\Sigma,\leq,\deleted{\unsure{f|\ell|\rho}}C)$, où $\Sigma$ est un ensemble, $\leq$ une relation d'ordre sur $\Sigma$, $C$ une application de $\Sigma$ en lui-même
\scanpage{51}
\mspage{4\deleted{\unsure{8}}7}
% archive sheet number: 50
satisfaisant les axiomes

\uline{Spat 1}. $\leq$ est une relation d'ordre, admettant un plus petit élément (noté $0_{\Sigma}$) et un plus grand élément (noté $1_{\Sigma}$), et \unsure{$=$|\textquotedbl} des Inf \inserted{quelconques (donc aussi des sup quelconques)}.

\uline{Spat 2} On \deleted{\unsure{est une|est un}} a $C^2=\mathrm{id}$ ($C$ est une involution) et \unsure{3)|b)} \deleted{\unsure{On a}}. $x\leq y\Longleftrightarrow Cy\leq Cx$ (i.e. $C$ est un \unsure{anti-automorphisme|antiautomorphisme} involutif de l'ens ordonné sur lui-même).

\uline{Spat \unsure{3)|3}} $\forall x$, on a $x\wedge Cx=0_{\Sigma}$ (ce qui équivaut, par \unsure{a)|2)}, à $\deleted{\unsure{x}}Cx\vee x=1_{\Sigma}$)

Dans un tel ensemble, on note
\[
\unsure{X|x}\mathrel{|\unsure{o|\circ}|}\unsure{Y|y}\Longleftrightarrow X\leq CY\quad(\Leftrightarrow Y\leq CX).
\]
\deleted{\unsure{On a}}
\[
X\mathrel{|o|}Y\Longrightarrow X\wedge\deleted{\unsure{C}}Y=0_{\Sigma}
\]
On trouve
\[
\begin{aligned}
X\mathrel{|o|}Y &\ \text{est relation symétrique}\\
X\mathrel{|o|}X &\Longleftrightarrow X=0_{\Sigma}
\end{aligned}
\]
\unsure{donc|alors} dans $\Sigma^{\unsure{*|\times}}=\Sigma\setminus\{0_{\Sigma}\}$, la relation $\mathrel{|o|}$ est \unsure{antiréflexive}, en plus d'être symétrique.

\uline{Ex} Soit $\unsure{\mathcal A|\mathscr A}$ un ens. muni de une relation $\mathrel{|o|}$ qui est symétrique, soit $\mathcal A_0$ l'ens des $X\in A$ tels que $X\mathrel{|o|}\deleted{\unsure{Y}}X$ \deleted{\unsure{pour tt}} \deleted{\unsure{On}} \deleted{\unsure{suppose que} $X\in\mathcal A_0$} \unsure{$Y$} ($\mathcal A_0=\unsure{\varnothing|\phi}$ dans le cas où $\mathcal A$ est un atelier, et $\mathrel{|o|}$ sa \deleted{\unsure{struc}} relation de disjonction intrins\unsure{...|èque}).

\inserted{On suppose que $X\in\mathcal A_0\Longrightarrow X\mathrel{|o|}Y\ \forall Y\in\mathcal A$.}

On trouve alors une notion de support et de cosupport, et un ensemble \deleted{\unsure{de}}
\[
\Sigma\subset\unsure{P|\mathcal P}(\mathcal A)
\]
\scanpage{52}
\mspage{48}
% archive sheet number: 51
\unsure{de|des} supports, stable dans $\unsure{P|\mathcal P}(\unsure{\mathcal A|\mathscr A})$ par $\bigcap\limits_i{}_{\unsure{s|\xi|\phi}}$ \unsure{quelconque|quelconques}, et donc ayant des Inf quelconques, donc aussi des Sup quelconques.
Le plus petit élément est $\mathcal A_{00}$\inserted{$=\{X\in\mathcal A\mid X\mathrel{|\unsure{o|\circ}|}Y\ \forall Y\in\mathcal A\unsure{)|\}}$}, le plus grand est $\mathcal A$. On a \inserted{une opération $C$ et} $C^2=\mathrm{id}$ (symétrie de $\mathrel{|o|}$). Ainsi, les conditions \deleted{a) b) c)}\inserted{Spat 1,2} d'un spatios sont satisfaites. On a \deleted{\unsure{d)}}\inserted{Spat 3} \unsure{sss|ssi} $X\mathrel{|o|}X\Longrightarrow X\in\mathcal A_{00}$ --- c'est la condition que j'avais mise dès le début, en douce.

Appliquons cette construction au cas $\mathcal A=\Sigma,\mathrel{|o|}_{\Sigma}$ ! On trouve bien $\Sigma_0=\{0_{\Sigma}\}=\Sigma_{00}$. Soit
% slanted
\margin{Si on part de $\mathcal A=\Sigma^{\unsure{*|\times}}=\Sigma\setminus\{0_{\Sigma}\}$, on \unsure{retrouve} encore $\Sigma$, comme ens des supports cette fois d'une relation symétrique et \uline{\unsure{antiréflexive}}. \unsure{Cas} \unsure{*} $\mathcal A=\unsure{\varnothing|\phi}$ sss $\Sigma$ réduit à un élément (\uline{Spatios} \uline{vide}).}
$A\unsure{\subset|\subseteq}\Sigma$, son cosupport est par définition
\[
\begin{aligned}
\operatorname{cosupp}(A)&=\{\deleted{\unsure{X}}Z\in\Sigma\mid\forall X\in A,\ \underbrace{X\mathrel{|o|}Z}_{\text{i.e. } Z\leq CX}\}\\
&=\{Z\in\Sigma,\ Z\leq\underset{X\in A}{\operatorname{Inf}}\,CX=\operatorname{Inf}CA\\
&\hspace{16em}=\deleted{\unsure{\operatorname{Inf}}}C\operatorname{Sup} A.\\[-2pt]
&=\widetilde{C\operatorname{Sup} A}
\end{aligned}
\]
\deleted{\unsure{Soit|Soient} \unsure{donc}}

Donc $\operatorname{cosupp}. A$ ne dépend que de $\operatorname{Sup} A$\inserted{$=\xi$, c'est $\widetilde{C(\xi)}$}. On aura de \unsure{$=$|même}
\[
\operatorname{supp} A\overset{\mathrm{def}}{=}\operatorname{cosupp}\operatorname{cosupp} A
\unsure{==|=}\underbrace{\widetilde{C\bigl(\underset{C(\xi)}{\tikz[baseline=(s.base)]{\node[inner sep=0pt] (s) {$\operatorname{Sup}\,\widetilde{C(\xi)}$};\draw ([yshift=-.2ex]s.south west)--([yshift=-.65ex]s.south west)--([yshift=-.65ex]s.south east)--([yshift=-.2ex]s.south east);}}\bigr)}}_{\xi}
=\widetilde{\xi}
\]
De cette façon, on voit que l'ensemble des supports s'identifie à $\Sigma$ via $\unsure{\Sigma|\xi}\longrightarrow\widetilde{\Sigma}$,
\scanpage{53}
\mspage{49}
% archive sheet number: 52
et via cette identification, le support d'une partie \inserted{$A$} de $\Sigma$ s'identifie à son Sup\inserted{$\xi$}, son cosupport à $C(\xi)$.

\deleted{Ceci dit, à un atelier \unsure{spatial} $(\unsure{\mathcal A|\mathscr A},\trianglelefteq,\mathrel{|\unsure{o|\circ}|})$ est associé un \unsure{spatios|Spatios} $(\Sigma,\leq,C)$, avec en plus une}

\unsure{Étant|Etant} donné un Spatios $\Sigma$, on \deleted{définit} appelle \uline{figure} \uline{\unsure{du|de}} \uline{Spatios}
\[
\deleted{\mathrm{Fig}\,\unsure{\mathrm{ind}|\mathrm{ord}}(\Sigma)}
\]
\[
\deleted{\mathrm{Fig}\,\mathrm{spat}(\Sigma)=\{(\deleted{\unsure{A|\mathcal A}}\unsure{\Phi|\varPhi},\rho)\mid\Phi\in}
\]
\unsure{un|une} \deleted{partie} couple $(\Phi,\rho)$, où
\[
\Phi\subset\Sigma
\]
est une partie de $\Sigma$, et $\rho$ une relation sur $\Phi$, satisfaisant les conditions

\uline{\deleted{\unsure{fig|Fig}}Fig spat 1}\quad $\rho$ est une relation d'ordre sur $\Phi$

\uline{Fig spat 2}\quad a) $\forall X\in\Phi$, on a $X\ne0_\Sigma$, et

\hspace*{4em}b) $\forall X,Y\in\Phi,\ X\ne Y\Longrightarrow X\mathrel{|o|}Y$ (i.e. $\unsure{x|X}\leq CY$)

\uline{NB} On pourrait être plus exigeant, en s'inspirant de l'axiome des supports \deleted{\unsure{de}}pour un atelier. On ne va pas s'en préoccuper pour le moment.

On désigne par $\mathrm{Fig}\,\mathrm{spat}(\deleted{\unsure{\mathcal A|A}}\Sigma)$ l'ensemble des figures spatiales de $\Sigma$, et par
\scanpage{54}
\mspage{50}
% archive sheet number: 53
$\mathrm{Fig}\,\mathrm{spat}\,\text{\unsure{éle|él|él.}}(\Sigma)$\inserted{ou $\mathrm{AtSpat}(\Sigma)$} l'ensemble des figures spatiales \uline{élémentaires} (i.e. telles que $\unsure{\Phi|\varPhi}$ ait un plus grand élément). On va mettre une structure d'atelier sur $\unsure{\mathrm{At}|\mathrm{AT}}(\Sigma)$, ainsi.
\[
\deleted{(\Phi,\rho)\trianglelefteq\Phi'}\qquad
\mathrm{AtSpat}(\Sigma)\subset\mathrm{Atens}(\Sigma)
\]
1) La relation $\trianglelefteq$ dans $\mathrm{Atspat}(\Sigma)$ est induite par la relation similaire dans $\mathrm{Atens}(\Sigma)$.

2) La relation $\mathrel{\overset{\unsure{\circ|o}}{<}\mkern-4mu<}$ dans $\mathrm{Atspat}\,\Sigma$ est \deleted{induite}\inserted{définie} comme plus haut : $(\Phi,\rho)\mathrel{\overset{\circ}{<}\mkern-4mu<}(\Phi',\rho')$ \unsure{sss|ssi} \deleted{par la relation similaire dans $\mathrm{Atens}\,\Sigma$} $\forall X\in\Phi,\ \exists X'=\psi(X)\in\Phi'$, avec $X\leq\deleted{\unsure{\psi}} \unsure{X'}\deleted{\unsure{(X)}}$ (NB. \deleted{(dans \unsure{Atens} pour $\unsure{\ll|\trianglelefteq}$, comme on voit} cet $X'$ est unique), et $\psi:\Phi\longrightarrow\Phi'$ est croissante \deleted{facilement \illegible{1} \unsure{constaté|constitué})} et transforme plus \unsure{grd|grand} \unsure{élémt|élément} en \unsure{idem|idm}.

3) La relation $\mathrel{|\unsure{o|\circ}|}$ dans $\mathrm{Atspat}(\Sigma)$ est \deleted{par contre} plus stricte que dans $\mathrm{Atens}(\Sigma)$ : on pose $(\Phi,\rho)\mathrel{|o|}(\Phi',\rho')$, sss (posant $X={}$plus grand \unsure{élémt|élément} de $\Phi$, et idem pour $\unsure{\alpha'|X'}$) non seulement on a $X\ne X'$ (relation $\mathrel{|o|}$ dans $\mathrm{Atens}$) mais même $X\mathrel{|o|}X'$ (ce qui est plus fort, vu que $X,\unsure{x'|X'}\ne0_\Sigma$).

Je dis que pour ces relations, on a \uline{un atelier}, \deleted{spatial} et même (je présume) un \uline{atelier} \uline{spatial}.
\scanpage{55}
\mspage{51}
% archive sheet number: 54
\uuline{\deleted{\unsure{Oblg|Obl|Fig}}At 1}\quad Clair

\uline{\deleted{\unsure{Ob|O}}At 2}\quad Clair \deleted{(c'est l'axiome ensembliste)}

\uline{At 3}\quad Clair

\uline{At 4}\quad Clair.

\deleted{\unsure{Cons|On}} Voyons quel est \unsure{le|l'} \deleted{\unsure{ens des}} \deleted{\unsure{supp}} spatios associé au précédent \unsure{Atelier|atelier}. \deleted{\unsure{Comme}} Soit $\unsure{\mathcal A|\mathscr A}=\mathrm{At}\,\mathrm{spat}(\Sigma)$. \unsure{Considérons|Consid.} \deleted{$(\unsure{\Phi|\varPhi},\rho)$} l'application
\[
\begin{aligned}
\mathcal A&\unsure{\text{\textemdash}|\longrightarrow}\Sigma\\
X=(\Phi,\rho)&\longmapsto\text{plus grand élément de }\Phi,\\
&\hspace{4em}\text{(noté }X^{\unsure{\circ|o}}\text{, cf justification \inserted{plus bas})}
\end{aligned}
\]
Ceci dit, la relation $\mathrel{|\unsure{o|\circ}|}$ dans $\mathcal A$ \unsure{est} induite, via cette application, par la relation de \unsure{$=$|même} nom dans $\Sigma$. On trouve donc une application injective
\[
\Sigma_{\mathcal A}\longrightarrow\Sigma_{\Sigma}
\underset{\inserted{\substack{\text{on vient de}\\\text{le voir}}}}{=}\Sigma.
\]
Mais on a aussi
\[
\begin{aligned}
\Sigma&\longrightarrow\mathcal A\\
\unsure{x|X}&\mathrel{\unsure{\vdash|\mapsto}}\deleted{\unsure{\{}}\bigl(\{x\}\unsure{,}\rho_x\deleted{\unsure{\}}}\bigr)
\qquad \rho_x\ \text{\unsure{la|l'} \inserted{\unsure{Unique|unique}} relation d'ordre sur }\{x\}
\end{aligned}
\]
telle que la relation $\mathrel{|o|}$ sur $\Sigma$ soit induite par celle sur $\mathcal A$, et le composé
\diagram{\begin{tikzcd}[ampersand replacement=\&,column sep=.32em]
\Sigma \& \unsure{\text{\textemdash}|\longrightarrow} \& \mathcal A \& \unsure{\text{\textemdash}|\longrightarrow} \& \Sigma
\end{tikzcd}}
est l'identité.
\scanpage{56}
\mspage{52}
% archive sheet number: 55
d'où
\[
\Sigma_{\unsure{\mathcal A|\mathscr A}}\simeq\Sigma_{\Sigma}\simeq\Sigma
\]
et on voit que ça respecte relations \unsure{d'ordres|d'ordre} et $C$.

Ceci dit, l'application canonique
\[
\mathcal A\longrightarrow\deleted{\unsure{\mathrm F}}\unsure{\mathrm{At}|\mathrm{Ft}}\,\mathrm{spat}(\Sigma_{\mathcal A})
\]
s'identifie, via $\Sigma_{\mathcal A}\simeq\Sigma$, en \unsure{composant}
\diagram{\begin{tikzcd}[ampersand replacement=\&,column sep=1.8em]
\mathcal A \arrow[r] \& \mathrm{Atspat}(\Sigma_{\mathcal A}) \arrow[r,"\sim"] \& \mathrm{Atspat}(\Sigma)=\mathcal A
\end{tikzcd}}
à l'application identique\unsure{,|.} Cette application canonique est donc non seulement injective, mais \uline{bijective} !

Revenons au cas d'un \unsure{Atelier|atelier} \deleted{spatial} quelconque $\mathcal A$. \deleted{\unsure{Il est}}\inserted{Posons $\Sigma=\Sigma_{\mathcal A}$.} Alors \deleted{$\mathcal A$ est \unsure{réalisé} comme sous-\unsure{Atelier|atelier} de $\mathrm{Atspat}(\Sigma)$}\inserted{on a un \uline{hom} \uline{\unsure{d'Ateliers|d'ateliers}}}
\[
\begin{aligned}
\mathcal A&\longrightarrow\mathrm{Atspat}(\Sigma)\\
X&\longmapsto\text{Dép}\,(\unsure{x|X})\qquad ,
\end{aligned}
\]
(même si \deleted{$\unsure{\Sigma|X}$} $\mathcal A$ non spatial). C'est un hom \uline{\unsure{fidèle|fidele}}, i.e.
\[
X\mathrel{|\unsure{o|\circ}|}Y\Longleftrightarrow\text{\unsure{Dep|Dép}}\,(x)\mathrel{|o|}\text{Dep}\,Y.
\]
Dire que $\mathcal A$ est spatial, signifie
\scanpage{57}
\mspage{53}
% archive sheet number: 56
qu'on a un hom injectif, qui identifie $\unsure{\mathcal A|\mathscr A}$ à un sous-atelier de $\mathrm{Atspat}(\Sigma)$, \deleted{\unsure{soit}} \inserted{noté \unsure{encore}} $\deleted{\unsure{\mathcal B|B}}\unsure{\mathcal A|A}$

Ainsi, on trouve, associé à $\mathcal A$, \unsure{$x$|:}

$1^{\circ}$) un Spatios $\Sigma$

$2^{\circ}$) un sous-atelier $\deleted{\unsure{\mathcal B|B}}\mathcal A\subset\mathrm{At}\,\mathrm{spat}(\Sigma)$. \deleted{\unsure{i.e.}}

Sous-atelier signifie ici, simplement, une \uline{partie}\inserted{de $\mathrm{Atspat}$}, qui \unsure{satisfasse|satisfait} la condition Sous at (cf p. 27) --- mais en fait, c'est \unsure{$=$|même} une partie \uline{fermée pour} $\trianglelefteq$.
Cette condition joue le rôle de la condition $\deleted{\unsure{\mathrm O|\mathrm{Ob}}}\mathrm{Atens}\ 1$ (cf \unsure{CF|CK} $\overline{\text{\uline{VIII}}}$\inserted{ p.~\unsure{309|39}}) pour les ateliers ensemblistes, rapportés à $L$\inserted{ (ens. des \unsure{lieux})} comme ens. de base \unsure{\}|?}. Il faut une condition qui $\underset{\text{-- -- -- --}}{\text{\deleted{\unsure{tienne}}}}$ lieu de $\mathrm{Atens}\ 2$ ($\forall x\in L,\ \{\{x\}\}\in A$). Ce sera Atspat 2. Ci dessous.

\deleted{Donc} On va donc rapporter la structure sur l'ens de base $\Sigma$, avec comme données $(\leq,\deleted{\unsure{\mathrel{|o|}|\trianglelefteq}}C,\mathcal A)$\deleted{\illegible{1}}, où $\leq,C$ font de $\Sigma$ un \unsure{spatios|Spatios}, et où $\mathcal A$ est un \deleted{\unsure{partie|part}} ens. \unsure{des|de} figures spatiales élémentaires \deleted{$\unsure{\mathcal A|A}$} relativement à $\Sigma$
\[
\mathcal A\subset\mathrm{At}\,\mathrm{spat}(\Sigma),
\]
